% !TEX root = ../Thesis.tex

\chapter{理论基础}\label{chap:theory}

{
本章将介绍论文研究工作所涉及的理论基础，包括多模光纤成像原理和深度学习神经网络的基本原理。

\section{多模光纤成像原理}\label{sec:mmf_principle}

\subsection{多模光纤的传输特性}

% TODO: 介绍多模光纤的基本结构、模式传输、模式色散等内容

\subsection{散斑的形成机制}

% TODO: 介绍光在多模光纤中传输后形成散斑的物理机制

\subsection{多模光纤成像的方法}

\subsubsection{基于波前整形的方法}

% TODO: 介绍基于波前整形的成像方法

\subsubsection{基于传输矩阵的方法}

% TODO: 介绍基于传输矩阵的成像方法

\subsubsection{基于深度学习的方法}

% TODO: 介绍基于深度学习的成像方法，引出本文的研究方向


\section{深度学习基础}\label{sec:dl_basics}

\subsection{卷积神经网络}

% TODO: 介绍CNN的基本原理、卷积层、池化层等

\subsection{图像分割网络}

% TODO: 介绍图像分割任务的特点

\subsection{UNet网络架构}

\subsubsection{UNet的网络结构}

% TODO: 详细介绍UNet的编码器-解码器结构、跳跃连接等

\subsubsection{UNet的优势}

% TODO: 分析UNet在图像重建任务中的优势

\subsubsection{UNet的变体}

% TODO: 介绍UNet的各种改进版本


\section{深度学习在光学成像中的应用}

% TODO: 综述深度学习在光学成像、散斑成像中的应用研究


\section{本章小结}

% TODO: 总结本章内容

}

